\section{Status quo: where the field stands and where the gap is}
\label{sec:related}

\plan{This section is written as real prose (not scaffold) because it is the
part of the paper that does not depend on our own unfinished data. Every
citation here was verified against Crossref, the arXiv API, or Europe PMC on
2026-08-04; the retrieval procedure is documented at the top of
\texttt{refs.bib}. Quantitative claims attributed to other groups are taken
verbatim from their abstracts.}

\subsection{Backbone generation is solved well enough that sequence design is
now the bottleneck}
\label{sec:related:backbones}

Deep generative models now produce enzyme backbones for a wide range of
chemistries. RFdiffusion established de novo design of protein structure and
function by denoising diffusion over backbone
coordinates~\cite{watson2023rfdiffusion}. RFdiffusion2 removed the requirement
that catalytic geometry be specified at the residue level: it designs enzymes
directly from functional-group geometries without specifying residue order or
performing inverse rotamer generation, and generates scaffolds for all 41
active sites in a diverse in-silico benchmark, compared with 16 for previous
methods~\cite{ahern2025rfdiffusion2}. Parallel efforts have produced designed
serine hydrolases~\cite{lauko2025serinehydrolases}, metallohydrolases with
designed metal sites~\cite{kim2025metallohydrolases}, and enzymes built by
catalytic motif scaffolding~\cite{braun2025motifscaffolding}.

\plan{The rhetorical work of this subsection is to concede generously. These
are strong papers and several share our senior author. The point is not that
backbone generation is inadequate --- it is that its success relocates the
bottleneck downstream, to sequence design and to what happens after the first
96 designs come back weakly active.}

Two features of this literature matter for what follows. First, the
sequence-design step is treated as a solved subroutine: backbones are passed to
ProteinMPNN~\cite{dauparas2022proteinmpnn} or
LigandMPNN~\cite{dauparas2025ligandmpnn} sampled at low temperature, and the
resulting sequences are filtered against structure-prediction and geometric
criteria. Second, the experimental readout is a small, one-shot test set:
RFdiffusion2 reports active candidates after testing fewer than 96 sequences
per catalytic mechanism~\cite{ahern2025rfdiffusion2}. Both features are
appropriate for demonstrating that a backbone generator works. Neither
addresses what to do when the resulting enzymes are, as they almost always
are, orders of magnitude less active than useful.

\subsection{Filtering is not search}
\label{sec:related:filtering}

The prevailing sequence-design protocol is generate-then-filter. Sequences are
sampled from a functionally blind inverse-folding prior, then scored against
structure-prediction confidence, backbone recapitulation, catalytic-geometry
and preorganization criteria --- AlphaFold3~\cite{abramson2024alphafold3},
Boltz~\cite{wohlwend2024boltz}, and ensemble methods such as PLACER, whose
introduction substantially improved experimental success rates in serine
hydrolase design~\cite{lauko2025serinehydrolases}.

\plan{This is the paper's central rhetorical move, so state it plainly and
early, and make sure the phrasing survives review: a filter is a selection
operator, not a search operator.}

A filter cannot create a sequence that passes; it can only surface the
sequences that a functionally blind prior happened to emit. This has two
consequences that the field has largely left implicit. Sampling temperature is
lowered to raise the pass rate, which means pass rate is purchased directly
with diversity: ProteinMPNN's own report notes that sequence diversity can be
increased considerably at higher temperature with only a small decrease in
recovery~\cite{dauparas2022proteinmpnn}, and the corresponding pass-rate cost
is why production pipelines sit near $T = 0.1$. And because the survivors are
drawn from a narrow neighborhood of the prior's mode, they are strongly
correlated with one another --- so a library's nominal size overstates the
number of genuinely independent attempts it represents.

\subsection{Reinforcement learning has reached inverse folding, but not
libraries}
\label{sec:related:rl}

Aligning generative models to objectives via reinforcement learning is now
standard practice, and Group Relative Policy Optimization
(GRPO)~\cite{shao2024deepseekmath} in particular has proven to be a
memory-efficient fit for the regime where rewards are computed by an external
oracle rather than learned. This machinery has begun to reach protein sequence
design. ProteinZero is an online RL framework for inverse folding that combines
structural guidance from ESMFold with a self-derived $\Delta\Delta G$
predictor, adds an embedding-level diversity regularizer to mitigate mode
collapse, and reports reducing design failure rates by 36--48\% relative to
ProteinMPNN, ESM-IF~\cite{hsu2022esmif}, and InstructPLM on the CATH-4.3
benchmark, with success rates above 90\% across diverse
folds~\cite{wang2025proteinzero}. Notably, its formulation balances three terms
--- multi-reward optimization, KL-divergence from a reference model, \emph{and}
diversity regularization --- so the anchor and the diversity mechanism are
separate components. Diversity-regularized direct preference
optimization has been applied to peptide inverse folding with related
motivation~\cite{park2024divdpo}.

\plan{ProteinZero is the closest prior art to our mechanism and must be
confronted in the first paragraph a reviewer reads on this topic, not buried.
The differentiation is real but it is about the \emph{objective}, not the
optimizer: (i) general folds on CATH vs.\ enzyme active sites with ligands and
catalytic geometry; (ii) sequence-level metrics --- designability, recovery,
stability --- vs.\ a library-level objective; (iii) a diversity
\emph{regularizer} defended as anti-mode-collapse hygiene vs.\ diversity as a
first-class design goal justified by screening economics; (iv) no physical
library. If compute allows, running ProteinZero as a baseline arm is the
strongest possible version of this argument. \todo{decide whether ProteinZero
is runnable as a baseline}}

A separate line of work has independently arrived at the structural move our
selection experiments test. Explorative Modeling factors the \emph{training
loop} rather than the generation procedure: it samples $K$ candidate matches
between model generations and data and trains on the best, so that predictions
commit to modes rather than blurring them, and reports that scaling exploration
acts as a third pretraining axis beyond parameters and data --- improving FLOP
efficiency \num{4.1}$\times$ and sample efficiency \num{6.2}$\times$, and
reaching \num{1.43} FID on ImageNet without
guidance~\cite{gladstone2026explorative}. Oversampling candidates and training
on a chosen subset is exactly the operation our E18 arms perform, and their
observation that the mode-seeking variant ``applies no pressure to cover every
mode, so on its own it is mode-seeking and can collapse'' is an independent
derivation of the failure we predict for pure quality-directed selection.

\plan{Cite this as convergent evidence, not as prior art we are extending, and
be precise about the disanalogy --- a reviewer who knows the paper will check.
Their selection criterion is the training objective itself (minimum
reconstruction loss against real data), so every candidate can be scored before
one is chosen. Ours cannot be: the objective requires a structure prediction,
which is the expensive step and the entire reason to oversample. We therefore
select on \emph{proxies} --- policy log-probability, sequence distance --- before
any reward is known. That is why a selection rule can lose to random in our
setting, which a minimum-loss rule cannot, and it is the reason E18 carries a
\texttt{random} control at matched oversample: their paper positions exploration
as a scaling axis where more candidates inherently help, and reports no ablation
separating selection quality from simply processing more samples.}

Read together, these two papers and our own results converge on the same
failure from three directions. ProteinZero introduces an embedding-level
diversity regularizer specifically to mitigate mode
collapse~\cite{wang2025proteinzero}; Explorative Modeling reports that its
mode-seeking variant applies no pressure to cover every mode and can collapse
on its own~\cite{gladstone2026explorative}; and we find that correcting the RL
objective converts a policy that diffuses outward into one that converges onto
a narrow solution, with positional entropy falling from ${\sim}1.7$ to
\num{1.06}. Collapse is not incidental to reward-aligned generative training ---
it is what the corrected objective does by default, and it is the reason a
library-level method cannot simply inherit a sequence-level one.

\plan{Where we differ is in what the anchor has to do. ProteinZero carries a
KL term \emph{and} a separate diversity regularizer; we carry only the KL term,
and E13 shows it is sufficient on its own --- sweeping its weight moves
per-design coverage $U_{15}$ from 510 to 881 monotonically, with total library
coverage peaking in an interior plateau. That is a claim about parsimony, not
about superiority: their regularizer operates in embedding space and ours does
not, and we have not run the two against each other. State it as ``one control
suffices in our setting'', which is defensible, rather than ``the regularizer
is unnecessary'', which we have not tested.}

What this line of work does not address is the library. Its objectives and
benchmarks are defined per sequence: is \emph{this} sequence designable, stable,
recovered. The quantity that determines an experimental outcome is a property
of the set --- how many distinct, genuinely different candidates clear the bar,
and how far apart they are.

\subsection{Library design assumes a natural enzyme}
\label{sec:related:libraries}

A mature literature does optimize libraries, but from the opposite starting
point. MODIFY learns from natural protein sequences to infer evolutionarily
plausible mutations and predict enzyme fitness, and explicitly co-optimizes
predicted fitness and sequence diversity of starting libraries, prioritizing
high-fitness variants while ensuring broad sequence
coverage~\cite{ding2024modify}. Active Learning-assisted Directed Evolution
(ALDE) uses uncertainty quantification over an epistatic five-residue active
site and improved the yield of a non-native cyclopropanation product from 12\%
to 93\% in three rounds of wet-lab experimentation~\cite{yang2025alde}.
Machine-learning-guided cell-free expression platforms scale the assay side,
mapping fitness across large variant sets~\cite{landwehr2025cellfree}.

\plan{MODIFY is the single most important citation for our diversity claim,
because it proves the field already accepts that fitness and diversity should
be co-optimized in a library. That is helpful --- we are not arguing for
diversity from scratch. What we are arguing is that every existing method for
doing so is unavailable in the de novo regime.}

Every one of these methods presupposes what a de novo enzyme lacks. MODIFY's
fitness prior is evolutionary, learned from natural homologs; a de novo
backbone has no homologs and no alignment. ALDE and active-learning methods
require a measurable starting fitness to bootstrap from; de novo designs
typically start at or near the detection limit. And both operate over a handful
of positions, whereas the mutations that rescue a de novo enzyme are, in our
hands, tens of substitutions from the starting design --- a combinatorial
volume that site-saturation and few-position libraries cannot reach.

\subsection{The gap}
\label{sec:related:gap}

Nobody has addressed library design for de novo enzymes. The regime is defined
by three simultaneous absences: no evolutionary prior, so MSA-based fitness
models do not apply; no starting activity, so active-learning loops have
nothing to condition on; and a target region of sequence space too far away for
local mutagenesis. The field's current answer --- low-temperature inverse
folding plus hard filters --- is not a search procedure at all, and it responds
to the difficulty of the regime by narrowing the search precisely when it
should be widening it.

\plan{Close on the thesis so the reader carries it into the results. Draft
formulation, to be tightened once F2 exists: function-aligned online RL
converts in-silico filters from a rejection step into a design objective,
producing de novo enzyme libraries that simultaneously pass the community's own
fidelity criteria at higher rates and span far more sequence space ---
increasing the number of independent experimental bets per unit of screening
effort.}

\todo{gap paragraph --- final phrasing of the thesis sentence, once the
headline figure is built}
